% SPDX-License-Identifier: Apache-2.0
% covers: Fifo
\datasheet{Fifo}{A FIFO of a power of two words, with a channel at each end.}

\dsfacts{\code{lib/parts/src/fifo.rs}}{\code{txhdl\_parts::fifo::Fifo}}%
{\code{//docs:runtime}, \code{//docs:examples}, \code{//docs:station}}

\subsection*{Function}

A channel is a buffer of two. \code{Fifo} is the depth a channel does
not have, for a unit whose consumer takes in bursts or holds off. It
has a channel in, \code{inp}, a channel out, \code{out}, and \code{D}
words between. \code{//docs:station} puts one behind a reservation
station for that reason.

\subsection*{Parameters}

\begin{center}\footnotesize
\begin{tabular}{@{}lp{0.7\textwidth}@{}}
\toprule
Parameter & Meaning \\
\midrule
\code{T} & What it holds: a \code{Transaction} and \code{Value} type. \\
\code{AW} & The width of its pointers, in bits. \\
\code{D} & How many words it holds, which must be \code{1 << AW}. \\
\bottomrule
\end{tabular}
\end{center}

\code{D} is stated rather than computed, because an expression in a
const parameter's position needs nightly Rust. The tables below use
\code{Fifo<U<8>, 2, 4>}, the FIFO of \code{ex\_queue}: four bytes.

\dstables{Fifo}

\subsection*{Behaviour}

\begin{itemize}
\item \code{mem} holds the words. \code{head} is where the oldest word
is, and \code{tail} is where the next word lands.
\item The pointers are equal both when it is empty and when it is
full. \code{full} says which.
\item A word is taken from \code{inp} whenever \code{full} is clear,
and written at \code{tail}, which moves on by one.
\item When it is not empty and \code{out} has room, the word at
\code{head} is sent on \code{out}, and \code{head} moves on by one.
\item A word taken lands in the memory at the end of its step and is
offered the step after. The source states the latency as one cycle.
\item \code{full} clears in a cycle that sends a word. It sets in a
cycle that takes a word and sends none, when the moved \code{tail}
meets \code{head}.
\end{itemize}

\subsection*{Verification}

\begin{itemize}
\item \code{fifo\_keeps\_every\_word\_in\_order} in
\code{txhdl\_parts}, in \code{//lib/parts:parts\_test}, runs
\code{Fifo<U<8>, 2, 4>} against a queue. It offers words at random
for 200 cycles and takes none for the first 40, so the FIFO fills. It
checks every word comes out in order, that \code{full} was seen, and
that it is empty at the end.
\item \code{ex\_queue} bursts nine words into the same FIFO while the
sink takes nothing for eight cycles, as \code{//docs:examples}
describes. The build simulates the netlist, named \code{queue4},
against that run under nvc and Verilator:
\code{//docs:queue\_sim\_queue4\_tb\_test} and
\code{//docs:queue\_vsim\_test}.
\end{itemize}

\subsection*{Used by}

The example \code{ex\_queue}. \code{ex\_fifo} is a different unit: an
asynchronous FIFO of its own, defined in the example.

\subsection*{Limits}

It runs on one clock, \code{DefaultClock}. \code{D} must be
\code{1 << AW}, and nothing checks it. When \code{full} is set it
takes no word, even in a cycle in which it sends one.
